% Options for packages loaded elsewhere
\PassOptionsToPackage{unicode}{hyperref}
\PassOptionsToPackage{hyphens}{url}
\documentclass[
]{article}
\usepackage{xcolor}
\usepackage[margin=1in]{geometry}
\usepackage{amsmath,amssymb}
\setcounter{secnumdepth}{-\maxdimen} % remove section numbering
\usepackage{iftex}
\ifPDFTeX
  \usepackage[T1]{fontenc}
  \usepackage[utf8]{inputenc}
  \usepackage{textcomp} % provide euro and other symbols
\else % if luatex or xetex
  \usepackage{unicode-math} % this also loads fontspec
  \defaultfontfeatures{Scale=MatchLowercase}
  \defaultfontfeatures[\rmfamily]{Ligatures=TeX,Scale=1}
\fi
\usepackage{lmodern}
\ifPDFTeX\else
  % xetex/luatex font selection
\fi
% Use upquote if available, for straight quotes in verbatim environments
\IfFileExists{upquote.sty}{\usepackage{upquote}}{}
\IfFileExists{microtype.sty}{% use microtype if available
  \usepackage[]{microtype}
  \UseMicrotypeSet[protrusion]{basicmath} % disable protrusion for tt fonts
}{}
\makeatletter
\@ifundefined{KOMAClassName}{% if non-KOMA class
  \IfFileExists{parskip.sty}{%
    \usepackage{parskip}
  }{% else
    \setlength{\parindent}{0pt}
    \setlength{\parskip}{6pt plus 2pt minus 1pt}}
}{% if KOMA class
  \KOMAoptions{parskip=half}}
\makeatother
\usepackage{color}
\usepackage{fancyvrb}
\newcommand{\VerbBar}{|}
\newcommand{\VERB}{\Verb[commandchars=\\\{\}]}
\DefineVerbatimEnvironment{Highlighting}{Verbatim}{commandchars=\\\{\}}
% Add ',fontsize=\small' for more characters per line
\usepackage{framed}
\definecolor{shadecolor}{RGB}{248,248,248}
\newenvironment{Shaded}{\begin{snugshade}}{\end{snugshade}}
\newcommand{\AlertTok}[1]{\textcolor[rgb]{0.94,0.16,0.16}{#1}}
\newcommand{\AnnotationTok}[1]{\textcolor[rgb]{0.56,0.35,0.01}{\textbf{\textit{#1}}}}
\newcommand{\AttributeTok}[1]{\textcolor[rgb]{0.13,0.29,0.53}{#1}}
\newcommand{\BaseNTok}[1]{\textcolor[rgb]{0.00,0.00,0.81}{#1}}
\newcommand{\BuiltInTok}[1]{#1}
\newcommand{\CharTok}[1]{\textcolor[rgb]{0.31,0.60,0.02}{#1}}
\newcommand{\CommentTok}[1]{\textcolor[rgb]{0.56,0.35,0.01}{\textit{#1}}}
\newcommand{\CommentVarTok}[1]{\textcolor[rgb]{0.56,0.35,0.01}{\textbf{\textit{#1}}}}
\newcommand{\ConstantTok}[1]{\textcolor[rgb]{0.56,0.35,0.01}{#1}}
\newcommand{\ControlFlowTok}[1]{\textcolor[rgb]{0.13,0.29,0.53}{\textbf{#1}}}
\newcommand{\DataTypeTok}[1]{\textcolor[rgb]{0.13,0.29,0.53}{#1}}
\newcommand{\DecValTok}[1]{\textcolor[rgb]{0.00,0.00,0.81}{#1}}
\newcommand{\DocumentationTok}[1]{\textcolor[rgb]{0.56,0.35,0.01}{\textbf{\textit{#1}}}}
\newcommand{\ErrorTok}[1]{\textcolor[rgb]{0.64,0.00,0.00}{\textbf{#1}}}
\newcommand{\ExtensionTok}[1]{#1}
\newcommand{\FloatTok}[1]{\textcolor[rgb]{0.00,0.00,0.81}{#1}}
\newcommand{\FunctionTok}[1]{\textcolor[rgb]{0.13,0.29,0.53}{\textbf{#1}}}
\newcommand{\ImportTok}[1]{#1}
\newcommand{\InformationTok}[1]{\textcolor[rgb]{0.56,0.35,0.01}{\textbf{\textit{#1}}}}
\newcommand{\KeywordTok}[1]{\textcolor[rgb]{0.13,0.29,0.53}{\textbf{#1}}}
\newcommand{\NormalTok}[1]{#1}
\newcommand{\OperatorTok}[1]{\textcolor[rgb]{0.81,0.36,0.00}{\textbf{#1}}}
\newcommand{\OtherTok}[1]{\textcolor[rgb]{0.56,0.35,0.01}{#1}}
\newcommand{\PreprocessorTok}[1]{\textcolor[rgb]{0.56,0.35,0.01}{\textit{#1}}}
\newcommand{\RegionMarkerTok}[1]{#1}
\newcommand{\SpecialCharTok}[1]{\textcolor[rgb]{0.81,0.36,0.00}{\textbf{#1}}}
\newcommand{\SpecialStringTok}[1]{\textcolor[rgb]{0.31,0.60,0.02}{#1}}
\newcommand{\StringTok}[1]{\textcolor[rgb]{0.31,0.60,0.02}{#1}}
\newcommand{\VariableTok}[1]{\textcolor[rgb]{0.00,0.00,0.00}{#1}}
\newcommand{\VerbatimStringTok}[1]{\textcolor[rgb]{0.31,0.60,0.02}{#1}}
\newcommand{\WarningTok}[1]{\textcolor[rgb]{0.56,0.35,0.01}{\textbf{\textit{#1}}}}
\usepackage{graphicx}
\makeatletter
\newsavebox\pandoc@box
\newcommand*\pandocbounded[1]{% scales image to fit in text height/width
  \sbox\pandoc@box{#1}%
  \Gscale@div\@tempa{\textheight}{\dimexpr\ht\pandoc@box+\dp\pandoc@box\relax}%
  \Gscale@div\@tempb{\linewidth}{\wd\pandoc@box}%
  \ifdim\@tempb\p@<\@tempa\p@\let\@tempa\@tempb\fi% select the smaller of both
  \ifdim\@tempa\p@<\p@\scalebox{\@tempa}{\usebox\pandoc@box}%
  \else\usebox{\pandoc@box}%
  \fi%
}
% Set default figure placement to htbp
\def\fps@figure{htbp}
\makeatother
\setlength{\emergencystretch}{3em} % prevent overfull lines
\providecommand{\tightlist}{%
  \setlength{\itemsep}{0pt}\setlength{\parskip}{0pt}}
\usepackage{bookmark}
\IfFileExists{xurl.sty}{\usepackage{xurl}}{} % add URL line breaks if available
\urlstyle{same}
\hypersetup{
  pdftitle={Codding Challenge 6},
  pdfauthor={Hikmat Paudel},
  hidelinks,
  pdfcreator={LaTeX via pandoc}}

\title{Codding Challenge 6}
\author{Hikmat Paudel}
\date{2026-03-26}

\begin{document}
\maketitle

{
\setcounter{tocdepth}{2}
\tableofcontents
}
\section{Question 1 (Regarding
Reproducibility)}\label{question-1-regarding-reproducibility}

The main purpose of creating your own functions and using iteration is
to keep your analysis consistent, reproducible, and efficient. Functions
let you reuse code instead of rewriting it, which helps minimize errors
and ensures the same approach is applied each time. Iteration tools like
loops automate repetitive tasks, making it easier to apply your analysis
to new or larger datasets.

\section{Question 2 (Conceptual)}\label{question-2-conceptual}

\begin{enumerate}
\def\labelenumi{\arabic{enumi})}
\item
  Function in R: A function in R is created by using the function() and
  assigning a name. Inputs are dfined in parentheses, and the code is
  written inside \{\}. It givess a result using return(). Functions are
  written in the scripts or R Markdown to resue code.
\item
  For loop in R: A for loop repeats code multiple times using for (i in
  \ldots), where i goes through a sequence. The code inside\{\} runs for
  each value of i. Results are usually stored in a variable that updates
  each iteration.
\end{enumerate}

\section{Question 3 (Load Data)}\label{question-3-load-data}

\begin{Shaded}
\begin{Highlighting}[]
\NormalTok{Cities }\OtherTok{\textless{}{-}}\FunctionTok{read.csv}\NormalTok{(}\StringTok{"Cities.csv"}\NormalTok{)}
\end{Highlighting}
\end{Shaded}

\section{Question 4 (Function
Writing)}\label{question-4-function-writing}

\begin{Shaded}
\begin{Highlighting}[]
\NormalTok{ Haversine}\OtherTok{\textless{}{-}} \ControlFlowTok{function}\NormalTok{(lat1, lon1, lat2, lon2) \{rad.lat1 }\OtherTok{\textless{}{-}}\NormalTok{ lat1 }\SpecialCharTok{*}\NormalTok{ pi}\SpecialCharTok{/}\DecValTok{180}
\NormalTok{  rad.lon1 }\OtherTok{\textless{}{-}}\NormalTok{ lon1 }\SpecialCharTok{*}\NormalTok{ pi}\SpecialCharTok{/}\DecValTok{180}
\NormalTok{  rad.lat2 }\OtherTok{\textless{}{-}}\NormalTok{ lat2 }\SpecialCharTok{*}\NormalTok{ pi}\SpecialCharTok{/}\DecValTok{180}
\NormalTok{  rad.lon2 }\OtherTok{\textless{}{-}}\NormalTok{ lon2 }\SpecialCharTok{*}\NormalTok{ pi}\SpecialCharTok{/}\DecValTok{180}
  
  \CommentTok{\# Haversine formula}
\NormalTok{  delta\_lat }\OtherTok{\textless{}{-}}\NormalTok{ rad.lat2 }\SpecialCharTok{{-}}\NormalTok{ rad.lat1}
\NormalTok{  delta\_lon }\OtherTok{\textless{}{-}}\NormalTok{ rad.lon2 }\SpecialCharTok{{-}}\NormalTok{ rad.lon1}
\NormalTok{  a }\OtherTok{\textless{}{-}} \FunctionTok{sin}\NormalTok{(delta\_lat }\SpecialCharTok{/} \DecValTok{2}\NormalTok{)}\SpecialCharTok{\^{}}\DecValTok{2} \SpecialCharTok{+} \FunctionTok{cos}\NormalTok{(rad.lat1) }\SpecialCharTok{*} \FunctionTok{cos}\NormalTok{(rad.lat2) }\SpecialCharTok{*} \FunctionTok{sin}\NormalTok{(delta\_lon }\SpecialCharTok{/} \DecValTok{2}\NormalTok{)}\SpecialCharTok{\^{}}\DecValTok{2}
\NormalTok{  c }\OtherTok{\textless{}{-}} \DecValTok{2} \SpecialCharTok{*} \FunctionTok{asin}\NormalTok{(}\FunctionTok{sqrt}\NormalTok{(a))}
  
  \CommentTok{\# Earth\textquotesingle{}s radius in kilometers}
\NormalTok{  earth\_radius }\OtherTok{\textless{}{-}} \DecValTok{6378137}
  
  \CommentTok{\# Calculate the distance}
\NormalTok{  distance\_km }\OtherTok{\textless{}{-}}\NormalTok{ (earth\_radius }\SpecialCharTok{*}\NormalTok{ c) }\SpecialCharTok{/} \DecValTok{1000}
  
  \FunctionTok{return}\NormalTok{(distance\_km) }
\NormalTok{\}}
\FunctionTok{colnames}\NormalTok{(Cities)}
\end{Highlighting}
\end{Shaded}

\begin{verbatim}
##  [1] "city"        "city_ascii"  "state_id"    "state_name"  "county_fips"
##  [6] "county_name" "lat"         "long"        "population"  "density"
\end{verbatim}

\section{Question 5}\label{question-5}

\begin{Shaded}
\begin{Highlighting}[]
\NormalTok{auburn }\OtherTok{\textless{}{-}}\NormalTok{ Cities}\SpecialCharTok{$}\NormalTok{lat[Cities}\SpecialCharTok{$}\NormalTok{city }\SpecialCharTok{==} \StringTok{"Auburn"}\NormalTok{ ]}
\NormalTok{auburn2 }\OtherTok{\textless{}{-}}\NormalTok{ Cities}\SpecialCharTok{$}\NormalTok{long[Cities}\SpecialCharTok{$}\NormalTok{city }\SpecialCharTok{==} \StringTok{"Auburn"}\NormalTok{ ]}

\NormalTok{nyc }\OtherTok{\textless{}{-}}\NormalTok{ Cities}\SpecialCharTok{$}\NormalTok{lat[Cities}\SpecialCharTok{$}\NormalTok{city }\SpecialCharTok{==} \StringTok{"New York"}\NormalTok{]}
\NormalTok{nyc2 }\OtherTok{\textless{}{-}}\NormalTok{ Cities}\SpecialCharTok{$}\NormalTok{long[Cities}\SpecialCharTok{$}\NormalTok{city }\SpecialCharTok{==} \StringTok{"New York"}\NormalTok{]}

\CommentTok{\# Calculate distance }
\NormalTok{ auburn\_nyc\_dist }\OtherTok{\textless{}{-}} \FunctionTok{Haversine}\NormalTok{(auburn, auburn2, nyc, nyc2)}

\NormalTok{auburn\_nyc\_dist}
\end{Highlighting}
\end{Shaded}

\begin{verbatim}
## [1] 1367.854
\end{verbatim}

\section{Question 6 (Loop)}\label{question-6-loop}

\begin{Shaded}
\begin{Highlighting}[]
\ControlFlowTok{for}\NormalTok{ (i }\ControlFlowTok{in} \DecValTok{1}\SpecialCharTok{:}\FunctionTok{nrow}\NormalTok{(Cities)) \{}
  
\NormalTok{  dist }\OtherTok{\textless{}{-}} \FunctionTok{Haversine}\NormalTok{(auburn, auburn2,}
\NormalTok{                    Cities}\SpecialCharTok{$}\NormalTok{lat[i],}
\NormalTok{                    Cities}\SpecialCharTok{$}\NormalTok{long[i])}
  
  \FunctionTok{print}\NormalTok{(dist)}
\NormalTok{\}}
\end{Highlighting}
\end{Shaded}

\begin{verbatim}
## [1] 1367.854
## [1] 3051.838
## [1] 1045.521
## [1] 916.4138
## [1] 993.0298
## [1] 1056.022
## [1] 1239.973
## [1] 162.5121
## [1] 1036.99
## [1] 1665.699
## [1] 2476.255
## [1] 1108.229
## [1] 3507.959
## [1] 3388.366
## [1] 2951.382
## [1] 1530.2
## [1] 591.1181
## [1] 1363.207
## [1] 1909.79
## [1] 1380.138
## [1] 2961.12
## [1] 2752.814
## [1] 1092.259
## [1] 796.7541
## [1] 3479.538
## [1] 1290.549
## [1] 3301.992
## [1] 1191.666
## [1] 608.2035
## [1] 2504.631
## [1] 3337.278
## [1] 800.1452
## [1] 1001.088
## [1] 732.5906
## [1] 1371.163
## [1] 1091.897
## [1] 1043.273
## [1] 851.3423
## [1] 1382.372
## [1] 0
\end{verbatim}

\section{Qusetion 7 (dataframe)}\label{qusetion-7-dataframe}

\begin{Shaded}
\begin{Highlighting}[]
\NormalTok{distance\_df }\OtherTok{\textless{}{-}} \ConstantTok{NULL}

\ControlFlowTok{for}\NormalTok{ (i }\ControlFlowTok{in} \DecValTok{1}\SpecialCharTok{:}\FunctionTok{nrow}\NormalTok{(Cities)) \{}
\NormalTok{    result\_i }\OtherTok{\textless{}{-}} \FunctionTok{data.frame}\NormalTok{(}
      \AttributeTok{City1 =}\NormalTok{ Cities}\SpecialCharTok{$}\NormalTok{city[i],}
      \AttributeTok{City2 =} \StringTok{"Auburn"}\NormalTok{,}
      \AttributeTok{Distance\_km =} \FunctionTok{Haversine}\NormalTok{(auburn, auburn2,}
\NormalTok{                              Cities}\SpecialCharTok{$}\NormalTok{lat[i],}
\NormalTok{                              Cities}\SpecialCharTok{$}\NormalTok{long[i]))}
    
\NormalTok{    distance\_df }\OtherTok{\textless{}{-}} \FunctionTok{rbind.data.frame}\NormalTok{(distance\_df, result\_i)}
\NormalTok{\}}
\FunctionTok{print}\NormalTok{(distance\_df)}
\end{Highlighting}
\end{Shaded}

\begin{verbatim}
##            City1  City2 Distance_km
## 1       New York Auburn   1367.8540
## 2    Los Angeles Auburn   3051.8382
## 3        Chicago Auburn   1045.5213
## 4          Miami Auburn    916.4138
## 5        Houston Auburn    993.0298
## 6         Dallas Auburn   1056.0217
## 7   Philadelphia Auburn   1239.9732
## 8        Atlanta Auburn    162.5121
## 9     Washington Auburn   1036.9900
## 10        Boston Auburn   1665.6985
## 11       Phoenix Auburn   2476.2552
## 12       Detroit Auburn   1108.2288
## 13       Seattle Auburn   3507.9589
## 14 San Francisco Auburn   3388.3656
## 15     San Diego Auburn   2951.3816
## 16   Minneapolis Auburn   1530.2000
## 17         Tampa Auburn    591.1181
## 18      Brooklyn Auburn   1363.2072
## 19        Denver Auburn   1909.7897
## 20        Queens Auburn   1380.1382
## 21     Riverside Auburn   2961.1199
## 22     Las Vegas Auburn   2752.8142
## 23     Baltimore Auburn   1092.2595
## 24     St. Louis Auburn    796.7541
## 25      Portland Auburn   3479.5376
## 26   San Antonio Auburn   1290.5492
## 27    Sacramento Auburn   3301.9923
## 28        Austin Auburn   1191.6657
## 29       Orlando Auburn    608.2035
## 30      San Juan Auburn   2504.6312
## 31      San Jose Auburn   3337.2781
## 32  Indianapolis Auburn    800.1452
## 33    Pittsburgh Auburn   1001.0879
## 34    Cincinnati Auburn    732.5906
## 35     Manhattan Auburn   1371.1633
## 36   Kansas City Auburn   1091.8970
## 37     Cleveland Auburn   1043.2727
## 38      Columbus Auburn    851.3423
## 39         Bronx Auburn   1382.3721
## 40        Auburn Auburn      0.0000
\end{verbatim}

\section{Question 8 (Github)}\label{question-8-github}

{[}CodingChallenge6{]}\url{https://github.com/Hikmat-au/Assignments-repository/tree/main/Coding\%20Challenge\%206}

\end{document}
